% /Users/pikachu/books/payments-book/src/parts/part5/ch29-processing-center.tex
\chapter{Процессинговый центр}
\label{ch:processing-center}
\margintoc


\begin{kaobox}[title=Что вы узнаете из этой главы]
  Что скрыто за \enquote{чёрным ящиком}, который за доли
  секунды отвечает на авторизационный запрос: из каких
  компонентов состоит процессинговый центр, как между ними
  делится бюджет задержки, почему схема active-active здесь
  обязательна, как процессинг встраивается в сети
  \scheme{Visa}, \scheme{Mastercard}, НСПК и СБП
  и какие требования к нему предъявляют регуляторы
  и стандарты безопасности.
\end{kaobox}

Процессинговый центр удобно мыслить как машину с тремя
инвариантами: ответить быстро, пережить отказ узла
и не выпустить ключевой материал за пределы защищённого
контура. Почти каждое архитектурное решение здесь приходится
оценивать по трём вопросам: что будет с задержкой, что будет
с доступностью и кто в итоге контролирует ключи.

\section{Компоненты процессингового центра}
\label{sec:pc-components}

Платёжный шлюз, описанный в предыдущей главе, принимает запрос
мерчанта и переводит его в сообщение \iso{8583}. Но на этом
работа не заканчивается. Между шлюзом и карточной сетью стоит
система, которая должна принять сообщение, выбрать маршрут,
выполнить криптографические проверки, записать след операции
и позже превратить авторизацию в клиринговую и расчётную
запись. Этой системой и является процессинговый центр.

\begin{kaobox}[title={Определение: Процессинговый центр}, colback=blue!4, colbacktitle=blue!15]
  Процессинговый центр (processing center) — комплекс
  программно-аппаратных средств, выполняющий авторизацию
  карточных транзакций, формирование клиринговых файлов,
  расчёт settlement-позиций и криптографические операции
  (верификация PIN, генерация и проверка криптограмм).
  В российском праве процессинговый центр может выступать
  как операционный и платёжный клиринговый центр (ОПКЦ),
  если имеет соответствующую лицензию ЦБ.
\end{kaobox}

Если смотреть на процессинг не как на один сервис, а как
на короткую производственную линию, в нём удобно выделить
четыре основных компонента. У каждого свой ритм работы,
свои точки отказа и свой способ ошибиться.

\subsection*{Authorization switch}

Authorization switch — это коммутатор авторизации, в который
попадает каждое входящее сообщение \iso{8583}. Он принимает
запросы от эквайеров или от шлюзов, действующих от их имени,
выбирает маршрут к нужному эмитенту или в карточную сеть
и возвращает ответ. Работает он поштучно и в реальном времени:
допустимая задержка здесь измеряется не секундами,
а десятками миллисекунд.

Маршрут внутри switch определяется BIN-таблицей:
для каждого BIN-диапазона известно, через какой канал
(VisaNet, Banknet, ОПКЦ НСПК, прямое подключение к эмитенту)
следует отправить запрос. BIN-таблица обновляется регулярно:
карточные сети публикуют изменения, и процессинг обязан
применять их в установленные сроки, иначе часть транзакций
начнёт уходить не туда и будет отклоняться без вины клиента
или мерчанта.

Помимо маршрутизации, switch выполняет базовую валидацию
сообщения: проверяет обязательные поля, формат MTI
и корректность bitmap, дополняет запрос данными, которые
требует конкретная сеть, и ведёт журнал транзакций.
Эта запись с точностью до миллисекунд нужна не только
для отладки, но и для последующей сверки с клиринговыми
файлами.

Когда внешний маршрут недоступен, включается резервная
авторизация (stand-in processing). Если карточная сеть
или эмитент не отвечают из-за таймаута или сетевого сбоя,
switch может принять решение сам, опираясь на локальные
правила: лимит суммы, BIN-диапазон, историю операций по карте.
Это всегда вынужденная мера. Процессинг временно берёт
на себя кредитный риск, который в обычном режиме несёт
эмитент, поэтому stand-in ограничивают малыми суммами
и заранее определёнными категориями транзакций.

\subsection*{HSM farm}

За криптографию в процессинговом центре отвечают
специализированные аппаратные модули HSM
(Hardware Security Module). Это физически защищённые
устройства, которые хранят ключи в среде, рассчитанной
на сопротивление вскрытию, и выполняют криптографические
операции внутри защищённого периметра
(см. главу \ref{ch:crypto}).

Пул HSM в процессинге решает несколько задач. При PIN-операции
switch извлекает PIN block из \de{52}, отправляет его
в HSM вместе с зональным ключом (Zone PIN Key), а HSM
переводит PIN block из формата эквайера в формат эмитента.
При EMV-транзакции тот же контур проверяет криптограмму ARQC:
получает данные операции и ключ эмитента или его производный
вариант, вычисляет ожидаемое значение и сравнивает его
с тем, что пришло от карты. Когда эмитент формирует ответ,
через HSM может быть создана и ARPC
(Authorization Response Cryptogram), которая возвращается
карте для взаимной аутентификации.

HSM — одно из главных узких мест по задержке. Каждая
криптографическая операция занимает порядка 1--5~мс,
и при нагрузке в тысячи транзакций в секунду одного
устройства не хватает. Поэтому процессинговые центры
используют HSM farm: пул из десятков устройств
с балансировкой нагрузки, где на каждом загружен одинаковый
набор ключей. Отказ одного HSM не должен останавливать
обработку: нагрузку подхватывают оставшиеся устройства,
а замена происходит без остановки сервиса.

\begin{kaobox}[title=Важно, colback=red!5, colbacktitle=red!20]
  HSM — единственный компонент процессинга, для которого
  горизонтальное масштабирование ограничено физически: каждое
  новое устройство требует церемонии загрузки ключей
  (key ceremony), которая проводится в защищённом помещении
  с участием нескольких хранителей ключей и занимает часы.
  Планирование
  ёмкости HSM farm должно учитывать не только текущую нагрузку,
  но и время, необходимое для ввода в эксплуатацию нового
  устройства.
\end{kaobox}

\subsection*{Clearing engine}

Если authorization switch живёт в режиме запроса-ответа,
то clearing engine живёт пакетами. Он запускается
по расписанию, обычно один-два раза в сутки, собирает
авторизованные транзакции за период и добавляет данные,
которых на этапе авторизации ещё не было: финальную сумму
после добавления чаевых (tip adjustment), результат адресной
проверки, параметры DCC-конверсии.

Затем engine формирует клиринговый файл в формате сети:
TC33 для \scheme{Visa}, IPM для \scheme{Mastercard}
или собственный формат НСПК (см. главу \ref{ch:reconciliation}).

На этом же этапе система сопоставляет каждую клиринговую
запись с авторизацией по ключевым полям: STAN, RRN, сумме
и дате. Так обнаруживаются расхождения: авторизация без
клиринга (orphan authorization), клиринг без авторизации
(late presentment, force post), несовпадение сумм
(partial capture, tip adjustment). Каждое такое расхождение
означает потенциальную потерю денег, поэтому clearing engine
формирует отчёт об исключениях, который операционная команда
разбирает вручную или по заранее заданным правилам.

\subsection*{Settlement module}

Settlement module считает итоговые нетто-позиции участников:
сколько каждый эмитент должен перечислить и сколько каждый
эквайер должен получить с учётом interchange, комиссий сети
и корректировок вроде chargeback, representment и fees.
Результатом становится набор проводок, который уходит
в расчётную инфраструктуру: для внутрироссийских транзакций
через НСПК и платёжную систему Банка России
(см. главу \ref{ch:banks}), для международных — через
расчётные банки карточных сетей.

Именно здесь деньги перестают быть обещанием в учётной
системе и действительно начинают двигаться между банками.
Поэтому ошибка в settlement — не тот технический сбой,
который можно безболезненно закрыть повтором, а прямая
финансовая потеря.

Даже эта короткая карта компонентов уже показывает,
что процессинг — не один сервер и не один протокол.
Это связка маршрутизации в реальном времени, пакетной
обработки, аппаратной криптографии и расчётного контура.
Следующий вопрос поэтому естественен: как уложить такую
систему в жёсткий бюджет задержки и одновременно сделать
её отказоустойчивой.

\section{Производительность и отказоустойчивость}
\label{sec:pc-performance}

Процессинговый центр — одна из немногих систем, где требования
к задержке и доступности не маркетинговое украшение,
а контрактное обязательство: карточные сети устанавливают SLA
на время ответа и штрафуют за его нарушение.

\subsection*{Latency budget}

Типичное требование карточной сети — ответ на авторизационный
запрос в течение 3--5~секунд от момента получения сообщения
до отправки ответа. Но 3--5~секунд — это максимум, включающий
время на доставку сообщения эмитенту и обратно; на долю
процессинга остаётся значительно меньше. Хорошо спроектированный
процессинговый центр укладывается в p99 $<$ 100 мс
на обработку одного авторизационного сообщения, и этот бюджет
распределяется примерно так:

\begin{itemize}
  \item Приём и разбор сообщения \iso{8583}: 1--2~мс.
  \item Поиск BIN и выбор маршрута: 1--3~мс.
  \item Обращение к HSM (PIN translation или проверка
        криптограммы): 3--10~мс.
  \item Запись в журнал транзакций: 1--2~мс.
  \item Сборка ответа и отправка: 1--2~мс.
  \item Сетевая задержка до карточной сети: 10--50~мс
        (зависит от географии).
\end{itemize}
Суммарно внутренняя обработка занимает 10--20~мс; основная
часть бюджета уходит на ожидание ответа от эмитента (через
карточную сеть), на которое процессинг повлиять не может.
Именно поэтому оптимизация внутренней обработки, хотя
и необходима, даёт ограниченный эффект: разница между
10 и 20~мс незаметна на фоне 500--2000~мс ожидания ответа
от удалённого эмитента.

\subsection*{Active-active и георезервирование}

Процессинговый центр, обрабатывающий карточные транзакции,
не может позволить себе плановый простой: даже несколько минут
недоступности — это тысячи отклонённых транзакций, недовольные
держатели карт и штрафы от карточных сетей. Поэтому стандартная
архитектура — active-active с двумя (или более) ЦОД,
каждый из которых способен обрабатывать полный объём трафика.

В active-active конфигурации трафик распределяется между ЦОД
в нормальном режиме (обычно 50/50 или по географическому
признаку), и при отказе одного весь трафик переключается
на оставшийся. Ключевая сложность — синхронизация состояния
между ЦОД. Журнал транзакций, BIN-таблицы, stand-in лимиты
и HSM-ключи должны быть согласованы в обоих центрах.

Для авторизации, которая по природе почти не хранит состояния
между запросами, синхронизация минимальна: достаточно, чтобы
BIN-таблицы и ключи были одинаковы, а журнал транзакций
реплицировался асинхронно. Для клиринга и settlement,
работающих с накопленным состоянием, требования строже:
потеря клиринговых данных означает потерю денег, и RPO
(Recovery Point Objective) должен быть близок к нулю.

\begin{kaobox}[title=Важно, colback=red!5, colbacktitle=red!20]
  Для авторизации RPO может составлять несколько секунд: потеря
  нескольких транзакций из журнала — неприятность,
  но не катастрофа, потому что авторизации можно восстановить
  из логов карточной сети. Для клиринга и settlement RPO
  должен быть нулевым: потерянная клиринговая запись —
  это транзакция, за которую эквайер не получит деньги
  или эмитент не спишет средства. Эта асимметрия определяет
  архитектуру репликации: авторизация допускает асинхронную
  репликацию ради минимальной латентности, а клиринг требует
  синхронной — ради консистентности.
\end{kaobox}

\begin{kaobox}[title=На практике]
  Шардинг по BIN-диапазонам — простейшая и наиболее
  распространённая стратегия горизонтального масштабирования
  коммутатора авторизации. BIN однозначно определяет эмитента
  и карточную сеть, а значит — маршрут и набор ключей;
  транзакции с разными BIN не зависят друг от друга
  и могут обрабатываться разными экземплярами switch
  без координации. При необходимости шард можно перенести
  на другой сервер или в другой ЦОД без влияния на остальные
  BIN-диапазоны.
\end{kaobox}

\subsection*{Аварийное восстановление}

Георезервирование, то есть размещение ЦОД в разных географических
точках, защищает уже не от отказа отдельного сервера,
а от региональных катастроф: пожара, наводнения, отключения
электричества в районе. Расстояние между центрами всегда
компромиссно. Чем дальше они друг от друга, тем лучше защита
от локальных событий, но тем выше задержка синхронной
репликации. Типичный диапазон — 50--200~км: этого обычно
достаточно, чтобы пережить локальный инцидент и при этом
удержать сетевую задержку ниже 5~мс.

RTO (Recovery Time Objective), то есть время восстановления
после полного отказа ЦОД, для процессинга измеряется секундами,
а не минутами. Переключение трафика на резервный центр должно
происходить автоматически, без ручного вмешательства,
по сигналу мониторинга. Для крупного процессингового центра
ожидаемый уровень — RTO $<$ 30~секунд.

Иначе говоря, отказоустойчивость в процессинге начинается
не с красивой схемы в презентации, а с очень прозаического
вопроса: что именно мы готовы потерять и сколько секунд
система имеет право молчать.

\section{Подключение к платёжным системам}
\label{sec:pc-connectivity}

После внутренней архитектуры возникает следующий вопрос:
к чему именно подключён процессинг и на каких правилах
держится это подключение. Процессинговый центр не живёт
в пустоте. Он встроен в сеть внешних контуров, и каждый
из них требует своих каналов связи, своей сертификации
и своей дисциплины обработки сообщений.

\subsection*{VisaNet}

Подключение к VisaNet — глобальной сети \scheme{Visa} —
осуществляется через выделенные каналы (leased lines
или VPN поверх интернета с шифрованием). Авторизация проходит
через BASE I — систему реального времени, обрабатывающую
авторизационные запросы; клиринг — через BASE II, принимающую
пакетные файлы. Процессинг, подключающийся к VisaNet, проходит
многоэтапную сертификацию: тестирование форматов сообщений,
криптографии, логики резервной авторизации, обработки reversal
и таймаутов \todo[inline]{Источник: если нужен точный состав
  приёмочных сценариев и последовательность сертификации, проверять
  Visa Acceptance Testing Guide для участников сети; в публичном
  доступе детальная матрица тестов не раскрыта.}.

\subsection*{Banknet}

\scheme{Mastercard} использует собственную сеть — Banknet,
архитектурно аналогичную VisaNet, но с рядом протокольных
отличий. Ключевое из них — DMSA (Dual Message System
Architecture), при которой \scheme{Mastercard} допускает как dual
message, когда авторизация и клиринг разделены, так
и single message, когда финансовая запись формируется
в одном сообщении, — в отличие от \scheme{Visa},
исторически ориентированной на dual message. Клиринговые
файлы \scheme{Mastercard}
передаются в формате IPM через систему
GCMS \todo[inline]{Источник: если нужен точный формат IPM и
  структура GCMS на уровне полей и пакетных потоков, проверять
  Mastercard Interface Processing Manual; в открытых материалах
  доступны только общие описания сети.}.

\subsection*{ОПКЦ НСПК}

Все внутрироссийские карточные транзакции — независимо
от того, выпущена карта \scheme{Visa}, \scheme{Mastercard}
или Мир, — маршрутизируются через операционный и платёжный
клиринговый центр (ОПКЦ) НСПК. Это означает, что процессинговый
центр российского банка подключается не напрямую к VisaNet
или Banknet для внутренних транзакций, а к ОПКЦ, который
выступает единым хабом: принимает авторизационные запросы
от эквайеров, маршрутизирует их к эмитентам и обеспечивает
клиринг и settlement (см. главу \ref{ch:mir}).

Подключение к ОПКЦ требует сертификации по стандартам НСПК
и поддержки специфических расширений \iso{8583}, используемых
в российской платёжной инфраструктуре — в частности, полей
для передачи данных платёжного приложения Мир (МРА)
с российской криптографией (ГОСТ).

\subsection*{СБП}

Подключение к СБП архитектурно отличается от карточного
процессинга: если карточные транзакции используют \iso{8583}
и двухфазную модель (авторизация + клиринг), то СБП работает
по API-протоколу с финальными платежами — нет авторизации
как отдельной фазы, нет chargeback, и расчёт происходит
в реальном времени (см. главу \ref{ch:sbp}). Процессинговый
центр, поддерживающий и карты, и СБП, вынужден совмещать
два принципиально разных потока обработки в одной
инфраструктуре.

\subsection*{Перегрузка и управляемая деградация}

Что происходит, когда карточная сеть не отвечает? Процессинговый
центр не может просто ждать: тысячи новых авторизаций продолжают
поступать, и если каждая из них повиснет в ожидании, очередь
быстро переполнится и парализует весь процессинг.

Грамотная реализация включает несколько уровней защиты.
Первый — таймаут на уровне соединения с сетью
(обычно 30--60~секунд), после которого транзакция либо отклоняется,
либо переводится в stand-in. Второй — автоматический
предохранитель (circuit breaker):
если доля таймаутов превышает порог (например, 50\% за последнюю
минуту), процессинг прекращает отправку в эту сеть и переключает
все транзакции на stand-in или альтернативный маршрут.
Третий — возврат сигнала перегрузки эквайерам: если процессинг перегружен,
он начинает отвечать \rc{96} (System Malfunction), сигнализируя
эквайерам о временной недоступности. Тогда они могут включить
собственную логику резервного переключения (failover).

Именно в такие моменты становится видно, что подключение
к сети — это не просто канал связи. Это ещё и набор правил,
которые не дают внешнему сбою утянуть за собой весь
внутренний контур.

\section{Лицензирование и сертификация}
\label{sec:pc-licensing}

Процессинговый центр работает на пересечении требований
регулятора, карточных сетей и индустриальных стандартов
безопасности, и каждый из этих уровней предъявляет свои
требования. Технически процессинг можно описывать как набор
серверов, HSM и каналов связи, но в реальном мире такой
контур существует только пока он одновременно легален,
сертифицирован и поддаётся аудиту.

\subsection*{Требования Банка России}

В российском праве процессинговый центр может выступать
как оператор услуг платёжной инфраструктуры (ОУПИ),
а также как операционный центр и платёжный клиринговый центр.
Для каждой роли предусмотрены свои лицензионные требования
в рамках \law{161-ФЗ}: минимальный капитал, требования
к бесперебойности, обязательства по предоставлению отчётности
в ЦБ (см. главу \ref{ch:regulation}).

Кроме того, процессинг обязан выполнять требования
\law{382-П} (защита информации при осуществлении переводов
денежных средств) и \law{683-П} (требования к защите
информации в кредитных организациях), которые устанавливают
конкретные технические меры — от шифрования каналов связи
до контроля доступа и логирования операций.

\subsection*{Сертификация у карточных сетей}

Каждая карточная сеть проводит собственную программу
сертификации для процессинговых центров, подключающихся
к её инфраструктуре. Сертификация включает:

\begin{itemize}
  \item Тестирование форматов сообщений: корректность заполнения
        элементов данных \iso{8583} для всех типов транзакций
        (авторизация, reversal, клиринг, chargeback).
  \item Тестирование криптографии: корректность PIN translation,
        проверки EMV-криптограмм, работы с ключами.
  \item Тестирование отказоустойчивости: обработка таймаутов,
        режима stand-in и обнаружения дубликатов.
  \item Проверку соответствия правилам сети:
        Core Rules для \scheme{Visa}, Network Rules
        для \scheme{Mastercard}.
\end{itemize}
Сертификация — не разовое событие: при существенных изменениях
в протоколах (новые MTI, новые элементы данных, изменения
в криптографии) карточные сети требуют ресертификации,
и процессинг обязан пройти её в установленные сроки.

\subsection*{PCI DSS}

Процессинговый центр по определению входит в CDE
(cardholder data environment): через него проходят PAN,
PIN block, CVV и криптограммы, то есть самые чувствительные
платёжные данные (см. главу \ref{ch:pci-dss}). Это означает
категорию PCI DSS Level 1 — максимальный контур, требующий
ежегодного аудита квалифицированным QSA
(Qualified Security Assessor) и квартальных ASV-сканов
(Approved Scanning Vendor).

Для процессинга, хранящего и обрабатывающего PIN, дополнительно
применяется PCI PIN Security — отдельный стандарт, регулирующий
физическую и логическую защиту PIN, HSM-процедуры,
управление ключами и их уничтожение.

Если собрать всё вместе, становится видно, что процессинговый
центр — это не просто технический посредник между шлюзом
и карточной сетью. Он одновременно держит скорость,
криптографическое доверие, денежную финальность и соответствие
внешним правилам. Именно поэтому вокруг него так много
операционной дисциплины: слишком много вещей ломается здесь
не громко, а дорого.

\begin{chaptersummary}
\begin{itemize}
  \item Процессинговый центр соединяет четыре слоя:
        коммутатор авторизации (authorization switch),
        пул HSM, клиринговый движок (clearing engine)
        и расчётный модуль (settlement module).
  \item Их общий инвариант прост: процессинг должен быстро
        маршрутизировать запросы, безопасно выполнять
        криптографию, собирать клиринг и считать итоговые
        расчётные позиции.
  \item Внутренняя обработка авторизации обычно укладывается
        в десятки миллисекунд; главная доля задержки живёт
        во внешнем пути до эмитента и обратно.
  \item Стандартная схема отказоустойчивости — active-active
        с двумя ЦОД: авторизация терпит асинхронную репликацию,
        а клиринг и settlement требуют почти нулевого RPO;
        ожидаемый RTO измеряется секундами, а не минутами.
  \item Подключение к VisaNet, Banknet, ОПКЦ НСПК и СБП —
        это не только транспорт, но и разные правила
        сертификации, форматов и деградации под отказом.
  \item Снаружи над процессингом стоят три слоя требований:
        правовая рамка Банка России, сертификация карточных
        сетей и контур PCI DSS Level 1 вместе
        с PCI PIN Security.
\end{itemize}
\end{chaptersummary}
